\section{Conclusion}

The popularity paradox is real: models trained on high-popularity datasets exhibit 21.21 percentage points larger generalization gaps than those trained on low-popularity alternatives. Popular benchmarks attract nearly 25$\times$ more optimization research, consistent with a co-evolution hypothesis. However, texture bias is not the mechanism---modern architectures show improved shape-based generalization.

We recommend complementing popular benchmarks with diverse alternatives. The benchmarks we trust most may be the ones we should trust least.
